\documentclass{fiche}
\metadonnees{9}{Le feu sous l'épinette}{Bloc 2 — Ce qui est possible : futur et modaux}

\begin{document}
\entete

\objectifs{%
\textbf{Langue} : \emph{must} et \emph{have to} ; \emph{should}, \emph{ought to},
\emph{had better} ; le regret (\emph{should have} + participe).\par
\textbf{Texte} : Jack London, \og To Build a Fire \fg (1908).\par
\textbf{Durée} : 1\,h\,30 — exercices 35 min, texte et questions 30 min, version 25 min.}

%% ===================================================================
\exercice[13 min]{Obligation et absence d'obligation}

\rappel{\textbf{Rappel.} \emph{Must} vient de celui qui parle, \emph{have to} d'une
contrainte extérieure. \emph{Must not} interdit ; \emph{need not} et \emph{don't have to}
dispensent.}

\consigne{Complétez avec \emph{must}, \emph{have to}, \emph{must not} ou \emph{need not}.
Plusieurs réponses sont parfois possibles.}

\begin{anglais}
\begin{enumerate}[label=\arabic*.,itemsep=2.2mm,leftmargin=*]
  \item At seventy-five below, a man \trou{must not} fail in his first attempt to build a
    fire.
  \item To light the fire he \trou{had to} remove his mittens.
  \item You \trou{need not} carry the tent; the others will take it.
  \item ``No man \trou{must} travel alone in the Klondike after fifty below,'' said the
    old-timer.
  \item Travellers \trou{have to} register at the post before leaving.
  \item You \trou{must not} touch the snow on the branches above the fire.
  \item We \trou{don't have to} start before dawn; the trail is short.
  \item He \trou{had to} look at his hand to know whether he was holding the twig.
  \item ``You \trou{must} keep your head,'' he told himself.
  \item I \trou{need not} swear to it, but I am sure I heard him.
  \item In such a cold, the wet feet \trou{must not} stay in the moccasins.
  \item She \trou{has to} stand the strain until the spring.
\end{enumerate}
\end{anglais}

\note{\emph{Must not} : 1, 6, 11. \emph{Have to} au passé (\emph{had to}) : 2, 8.
\emph{Need not} / \emph{don't have to} : 3, 7, 10.}

\note[Deux pièges]{\emph{Must} n'a pas de passé : on emprunte \emph{had to} (2, 8).
Et \emph{must not} n'est pas la négation de \emph{must} : \emph{you must not go} interdit,
\emph{you need not go} dispense. Le français \og il ne faut pas \fg{} correspond à
\emph{must not} ; \og il n'est pas nécessaire \fg{} à \emph{need not}.}

%% ===================================================================
\exercice[10 min]{Conseil et regret}

\rappel{\textbf{Rappel.} \emph{Should} et \emph{ought to} conseillent ; \emph{had better}
avertit d'un danger immédiat. \emph{Should have} + participe dit ce qu'il aurait fallu faire,
et qu'on n'a pas fait.}

\consigne{Complétez.}

\begin{anglais}
\begin{enumerate}[label=\arabic*.,itemsep=2.2mm,leftmargin=*]
  \item He \trou{should not have built} the fire under the spruce tree.
  \item He \trou{should have built} it in the open.
  \item You \trou{had better} start now, or the light will go.
  \item Travellers \trou{ought to} listen to the old-timers.
  \item I \trou{should have listened} to the advice he gave me last fall.
  \item You \trou{had better not} take off your mittens in this wind.
  \item We \trou{should not have} come this far alone.
  \item He \trou{ought to have taken} a trail-mate with him.
  \item You look exhausted; you \trou{should} rest before the next stage.
  \item They \trou{should have} reached the camp by six. \emph{(supposition)}
\end{enumerate}
\end{anglais}

\note[Deux valeurs de \emph{should}]{Phrase 10 : \emph{should have} n'exprime pas toujours le
regret ; il sert aussi à la supposition raisonnable (\og ils devraient être arrivés \fg).
Seul le contexte tranche. \emph{Had better} est plus pressant que \emph{should} : il sous-entend
une conséquence fâcheuse, et sa négation est \emph{had better not}, jamais \emph{had not
better}.}

%% ===================================================================
\exercice[12 min]{Le froid et le feu}
\consigne{a) Associez chaque mot à sa traduction.}

\begin{multicols}{2}
\begin{enumerate}[label=\arabic*.,itemsep=1.2mm,leftmargin=*]
  \item a twig \hfill \trou{c}
  \item a bough \hfill \trou{h}
  \item a spruce \hfill \trou{a}
  \item the brush \hfill \trou{k}
  \item to squat \hfill \trou{e}
  \item a mitten \hfill \trou{b}
  \item to freeze \hfill \trou{j}
  \item the frost \hfill \trou{f}
  \item to melt \hfill \trou{d}
  \item to crackle \hfill \trou{l}
  \item the trail \hfill \trou{i}
  \item foot-gear \hfill \trou{g}
\end{enumerate}
\columnbreak
\begin{enumerate}[label=\alph*.,itemsep=1.2mm,leftmargin=*]
  \item une épinette, un sapin
  \item une moufle
  \item une brindille
  \item fondre
  \item s'accroupir
  \item le gel, le givre
  \item les chaussures
  \item une grosse branche
  \item la piste
  \item geler
  \item les broussailles
  \item crépiter
\end{enumerate}
\end{multicols}

\note[Une série]{\emph{To freeze} (freeze, froze, frozen) donne \emph{the frost} le gel,
\emph{frosty} glacial, \emph{frozen} gelé. Même alternance de voyelle que \emph{speak, spoke,
spoken} ou \emph{steal, stole, stolen}.}

\consigne{b) Complétez avec un mot de la liste.}
\begin{anglais}
\begin{enumerate}[label=\arabic*.,itemsep=2mm,leftmargin=*]
  \item He \trou{squatted} in the snow and fed the flame with dry \trou{twigs}.
  \item He had been forced to remove his \trou{mittens}, and his fingers had gone numb.
  \item Each \trou{bough} of the \trou{spruce} carried a heavy weight of snow.
  \item His wet \trou{foot-gear} would \trou{freeze} as soon as he stopped walking.
\end{enumerate}
\end{anglais}

%% ===================================================================
\newpage
\section{Texte}

\noindent\small\textit{Un homme marche seul sur la piste du Klondike par soixante-quinze
degrés au-dessous de zéro, avec un chien pour seule compagnie. Il a traversé une plaque de
glace : ses pieds sont trempés. Il vient d'allumer un feu sous un sapin.}\normalsize

\begin{texte}
He knew there must be no failure. When it is seventy-five below zero, a man must not fail in
his first attempt to build a fire --- that is, if his feet are wet. If his feet are dry, and
he fails, he can run along the trail for half a mile and restore his circulation. But the
circulation of wet and freezing feet cannot be restored by running when it is seventy-five
below. No matter how fast he runs, the wet feet will freeze the harder.

All this the man knew. The \gl[an old-timer]{old-timer}{un ancien, un vieux de la vieille} on
Sulphur Creek had told him about it the previous fall, and now he was appreciating the
advice. Already all sensation had gone out of his feet. To build the fire he had been forced
to remove his \gl[a mitten]{mittens}{une moufle}, and the fingers had quickly gone numb. His
\gl[a pace]{pace}{une allure, un pas} of four miles an hour had kept his heart pumping blood
to the surface of his body and to all the extremities. But the instant he stopped, the action
of the pump eased down. The cold of space \gl[to smite, smote]{smote}{frapper} the unprotected
tip of the planet, and he, being on that unprotected tip, received the full force of the blow.
The blood of his body \gl[to recoil]{recoiled}{reculer} before it. The blood was alive, like
the dog, and like the dog it wanted to hide away and cover itself up from the fearful cold. So
long as he walked four miles an hour, he pumped that blood,
\gl[willy-nilly]{willy-nilly}{bon gré mal gré}, to the surface; but now it
\gl[to ebb]{ebbed}{refluer} away and sank down into the \gl[a recess]{recesses}{un
repli, un recoin} of his body. The extremities were the first to feel its absence. His wet
feet froze the faster, and his exposed fingers numbed the faster, though they had not yet
begun to freeze. Nose and cheeks were already freezing, while the skin of all his body
chilled as it lost its blood.

But he was safe. Toes and nose and cheeks would be only touched by the frost, for the fire
was beginning to burn with strength. He was feeding it with twigs the size of his finger. In
another minute he would be able to feed it with branches the size of his wrist, and then he
could remove his wet foot-gear, and, while it dried, he could keep his naked feet warm by the
fire, rubbing them at first, of course, with snow. The fire was a success. He was safe. He
remembered the advice of the old-timer on Sulphur Creek, and smiled. The old-timer had been
very serious in laying down the law that no man must travel alone in the Klondike after fifty
below. Well, here he was; he had had the accident; he was alone; and he had saved himself.
Those old-timers were rather \gl[womanish]{womanish}{efféminé, de bonne femme --- c'est
l'homme qui pense, non le narrateur}, some of them, he thought. All a man had to do was to
keep his head, and he was all right. Any man who was a man could travel alone. But it was
surprising, the rapidity with which his cheeks and nose were freezing. And he had not thought
his fingers could go lifeless in so short a time. Lifeless they were, for he could scarcely
make them move together to grip a twig, and they seemed remote from his body and from him.
When he touched a twig, he had to look and see whether or not he had hold of it. The wires
were pretty well down between him and his finger-ends.

All of which counted for little. There was the fire, snapping and crackling and promising life
with every dancing flame. He started to untie his \gl[a moccasin]{moccasins}{un mocassin}.
They were coated with ice; the thick German socks were like \gl[a sheath]{sheaths}{une gaine,
un fourreau} of iron half-way to the knees; and the moccasin strings were like rods of steel
all twisted and knotted as by some \gl[a conflagration]{conflagration}{un incendie}. For a
moment he \gl[to tug]{tugged}{tirer d'un coup sec} with his numbed fingers, then, realizing
the folly of it, he drew his \gl[a sheath-knife]{sheath-knife}{un couteau de chasse}.
\end{texte}
\source{Jack London, \og To Build a Fire \fg, \emph{Lost Face}, New York, Macmillan, 1910.}

\partiel[15 min]{Questions}
\consigne{\en{Answer in English, in full sentences.}}

\begin{enumerate}[label=\arabic*.,itemsep=2mm,leftmargin=*]
  \item \en{Why must this first fire succeed?}
    \lignes{3}
    \corr{Because his feet are wet. A man with dry feet who fails can run half a mile and
    restore his circulation, but wet freezing feet cannot be warmed by running at
    seventy-five below: the faster he runs, the harder they freeze.}
  \item \en{What happens to his blood when he stops walking?}
    \lignes{3}
    \corr{Walking kept his heart pumping blood to the surface and the extremities. The
    instant he stops, the blood recoils from the cold and sinks into the recesses of his
    body, so the feet, fingers, nose and cheeks are the first to freeze.}
  \item \en{What did the old-timer tell him, and what does he think of that advice now?}
    \lignes{3}
    \corr{That no man must travel alone in the Klondike after fifty below. Now that he has
    had the accident and saved himself alone, he decides the old-timers are rather womanish,
    and that all a man has to do is keep his head.}
  \item \en{Find the details that contradict him at the very moment he congratulates
    himself.}
    \lignes{3}
    \corr{His cheeks and nose are freezing surprisingly fast; his fingers are already
    lifeless and can scarcely grip a twig; he has to look to know whether he is holding
    something.}
  \item \en{``The wires were pretty well down between him and his finger-ends.'' Explain the
    image.}
    \lignes{2}
    \corr{The nerves are compared to telegraph wires that have come down in a storm: the
    hands are still there but no message reaches them, and the body is becoming a country
    cut off from itself.}
  \item \en{The dog is mentioned twice. What is the comparison worth?}
    \lignes{2}
    \corr{The blood, like the dog, obeys an instinct older than thought: both want to hide
    from the cold. The animal and the man's own body know what his reasoning refuses to
    admit.}
\end{enumerate}

\note[Les modaux du texte]{\emph{There must be no failure}, \emph{a man must not fail},
\emph{no man must travel alone} : l'obligation est présentée comme une loi de la nature, non
comme un ordre. \emph{He can run}, \emph{cannot be restored}, \emph{he could remove},
\emph{he would be able to} : la capacité, qui rétrécit de ligne en ligne. \emph{All a man had
to do was to keep his head} : c'est l'homme qui se rassure, et le lecteur entend déjà que
c'est faux.}

%% ===================================================================
\newpage
\section{Version}
\consigne{Traduisez le passage en français. Il suit immédiatement le texte.}

\begin{version}
But before he could cut the strings, it happened. It was his own fault or, rather, his
mistake. He should not have built the fire under the spruce tree. He should have built it in
the open. But it had been easier to pull the twigs from the brush and drop them directly on
the fire. Now the tree under which he had done this carried a weight of snow on its boughs. No
wind had blown for weeks, and each bough was fully \gl[to freight]{freighted}{charger}. Each
time he had pulled a twig he had communicated a slight agitation to the tree --- an
imperceptible agitation, so far as he was concerned, but an agitation sufficient to bring
about the disaster. High up in the tree one bough \gl[to capsize]{capsized}{renverser,
faire chavirer} its load of snow. This fell on the boughs beneath, capsizing them. This
process continued, spreading out and involving the whole tree. It grew like an avalanche, and
it descended without warning upon the man and the fire, and the fire was
\gl[to blot out]{blotted out}{effacer, anéantir}! Where it had burned was a
\gl[a mantle]{mantle}{un manteau, une couche} of fresh and disordered snow.

The man was shocked. It was as though he had just heard his own sentence of death. For a
moment he sat and stared at the spot where the fire had been. Then he grew very calm. Perhaps
the old-timer on Sulphur Creek was right. If he had only had a
\gl[a trail-mate]{trail-mate}{un compagnon de route} he would have been in no danger now. The
trail-mate could have built the fire.
\end{version}
\source{\emph{Ibid.}}

\lignes{22}

\corr{\textbf{Proposition de traduction.} Mais avant qu'il eût pu couper les lacets, cela
arriva. C'était sa faute, ou plutôt son erreur. Il n'aurait pas dû faire son feu sous le
sapin. Il aurait dû le faire à découvert. Seulement il était plus commode d'arracher les
brindilles aux broussailles et de les laisser tomber directement dans le feu. Or l'arbre sous
lequel il avait fait cela portait sur ses branches un poids de neige. Le vent n'avait pas
soufflé depuis des semaines, et chaque branche était chargée à plein. Chaque fois qu'il avait
tiré une brindille, il avait communiqué à l'arbre une légère secousse --- imperceptible pour
lui, mais suffisante à provoquer le désastre. Tout en haut de l'arbre, une branche renversa
son chargement de neige. Celui-ci tomba sur les branches du dessous et les renversa à leur
tour. Le processus se poursuivit, s'étendant et gagnant l'arbre entier. Il grossit comme une
avalanche et s'abattit sans avertissement sur l'homme et sur le feu, et le feu fut anéanti !
Là où il avait brûlé s'étendait un manteau de neige fraîche et bouleversée.

L'homme resta saisi. C'était comme s'il venait d'entendre prononcer sa propre sentence de
mort. Un instant il demeura assis à fixer l'endroit où le feu avait été. Puis il devint très
calme. Peut-être l'ancien de Sulphur Creek avait-il raison. S'il avait seulement eu un
compagnon de route, il ne serait en aucun danger à présent. Le compagnon, lui, aurait pu
faire le feu.}

\note[Choix de traduction 1]{\emph{He should not have built\ldots He should have built\ldots} :
le regret, à rendre par \og il n'aurait pas dû \fg{} / \og il aurait dû \fg. Ne pas confondre
avec \emph{he must have built}, \og il a sûrement fait \fg{} (fiche 8).}

\note[Choix de traduction 2]{\emph{If he had only had a trail-mate he would have been in no
danger now} : irréel du passé dans la subordonnée, mais \emph{now} dans la principale ---
d'où en français un conditionnel présent, non passé : \og il ne serait pas en danger à
présent \fg. C'est le type mixte des conditionnelles, revu en fiche 16.}

\note[Choix de traduction 3]{\emph{so far as he was concerned} : \og en ce qui le
concerne \fg, ici \og pour lui \fg. Formule figée à reconnaître, comme \emph{as far as I
know}.}

\note[Choix de traduction 4]{\emph{Where it had burned was a mantle of fresh snow} : la
subordonnée est sujet du verbe, et l'ordre est inversé. Le français rétablit : \og Là où il
avait brûlé s'étendait un manteau\ldots \fg, avec inversion elle aussi, ce qui garde
l'effet.}

\note[Choix de traduction 5]{\emph{The trail-mate could have built the fire} : \emph{could
have} + participe, irréel du passé, \og aurait pu \fg. Ajouter \og lui \fg{} en français
rend le contraste que l'anglais obtient par la seule place du mot.}

%% ===================================================================
\partiel{Lexique à mémoriser}
\begin{itemize}[label=--,itemsep=0.8mm,leftmargin=*]\small
  \item \textbf{a failure} — un échec ; \textit{verbe} to fail \textit{(fiche 3)}
  \item \textbf{an attempt} — une tentative ; \textit{verbe} to attempt
  \item \textbf{the trail} — la piste ; a trail-mate \textit{« un compagnon de route »}
  \item \textbf{a pace} — une allure ; to keep pace with
  \item \textbf{a blow} — un coup ; \textit{verbe} to blow \textit{« souffler »}
  \item \textbf{to chill} — refroidir ; \textit{nom} a chill \textit{« un coup de froid »}
  \item \textbf{a twig} — une brindille
  \item \textbf{a bough} — une grosse branche ; \textit{rime avec} now
  \item \textbf{the brush} — les broussailles ; \textit{et} une brosse
  \item \textbf{to squat} — s'accroupir
  \item \textbf{to grip} — saisir, empoigner ; \textit{nom} a grip
  \item \textbf{to untie} — dénouer ; un- \textit{défait l'action :} to untie, to undo, to unlock
  \item \textbf{a rod} — une tige, une barre
  \item \textbf{folly} — la folie, la sottise ; \textit{adjectif} foolish
  \item \textbf{fault} — la faute (la responsabilité) ; \textit{à distinguer de} a mistake \textit{(l'erreur)}
  \item \textbf{in the open} — à découvert, en plein air
  \item \textbf{beneath} — au-dessous de ; \textit{plus littéraire que} under
  \item \textbf{a warning} — un avertissement ; \textit{verbe} to warn
  \item \textbf{a sentence} — une sentence, une condamnation ; \textit{et} une phrase
  \item \textbf{to grow} — devenir ; grew, grown \textit{— devant un adjectif, « devenir » et non « pousser »}
\end{itemize}

\end{document}
